\begin{slide}
    \begin{block}{Problem}
        A tree of $n$ nodes and starting positions (nodes) of a squirrel and two ravens are given. On a signal of Odin, one raven flies into the air and lands again. Meanwhile, the squirrel can travel in the tree, but may not pass over a node occupied by the other raven.
        
        Output the minimum number of signals after which the ravens are guaranteed to have captured the squirrel.
    \end{block}
    \begin{block}{Possible Approaches}
        \begin{enumerate}
        \item Dynamic programming on possible states
        \item Force the squirrel into the shallowest subtree
        \item Find the `highest' node the squirrel can reach
        \end{enumerate}
    \end{block}
\end{slide}

\begin{slide}
%    \begin{block}{Problem}
%        A tree and starting positions (nodes) of a squirrel and two ravens are given. On a signal, one raven flies into the air and lands again. Meanwhile, the squirrel can travel in the tree, but may not pass over a node occupied by the other raven.
%        
%        Output the minimum number of signals after which the ravens are guaranteed to have captured the squirrel.
%    \end{block}
    \begin{block}{Approach 1: Dynamic Programming}
        \begin{itemize}
        \item Idea: Dynamic programming by backwards breadth-first-search on the states
        \item The final states are when the squirrel is at the same position as a raven
        \item To move between states, consider the possible ways for the ravens and the squirrel to travel (bearing in mind the squirrel cannot move past a raven)
        \item Already computed states must be skipped to avoid an infinite loop
        \item Runtime $O(n^4)$
        \end{itemize}
    \end{block}
\end{slide}

\begin{slide}
%    \begin{block}{Problem}
%        A tree and starting positions (nodes) of a squirrel and two ravens are given. On a signal, one raven flies into the air and lands again. Meanwhile, the squirrel can travel in the tree, but may not pass over a node occupied by the other raven.
%        
%        Output the minimum number of signals after which the ravens are guaranteed to have captured the squirrel.
%    \end{block}
    \begin{block}{Approach 2: DFSs on the tree}
        \begin{itemize}
        \item Idea: Drive the squirrel into the shallowest subtree
        \item For every node, root the tree at this node and compute its depth
        \item The minimum depth is the number of signals required
        \item Runtime: $O(n^2)$
        \end{itemize}
    \end{block}
\end{slide}


\begin{slide}
%    \begin{block}{Problem}
%        A tree and starting positions (nodes) of a squirrel and two ravens are given. On a signal, one raven flies into the air and lands again. Meanwhile, the squirrel can travel in the tree, but may not pass over a node occupied by the other raven.
%        
%        Output the minimum number of signals after which the ravens are guaranteed to have captured the squirrel.
%    \end{block}
    \begin{block}{Approach 3: Assign heights to nodes}
        \begin{itemize}
        \item Idea: Node heights correspond to the number of signals required (minus $1$)
        \item Leaves are assigned a height of $0$
        \item A node has a height of $k$ if it becomes a leaf after all nodes of height $< k$ have been removed
        \item The squirrel should travel to the highest node (a `center' of the tree) it can reach
        \item Compute the highest position the squirrel can reach from its starting position
        \item If there are two centers, add $1$ for an extra signal
        \item Runtime: $O(n)$
        \end{itemize}
    \end{block}
\end{slide}